
\uline{factorisation de} $f$ \uline{comme composé d'une immersion et d'un morphisme lisse}, \uline{le complexe} $\underline{L}^{X/Y}_\bullet$ \uline{défini en} 2.1 \uline{est strictement parfait}.

Soit $f=g\circ i$, où $g$ est un morphisme lisse, et $i$ une immersion. D'après 1.2, $i$ est une immersion régulière. Si $\underline{J}$ est l'Idéal définissant $i$, $\underline{J}/\underline{J}^2$ est donc un $\underline{O}_X$-Module localement libre de type fini. De même, comme $g:V\to Y$ est lisse, $i^*(\underline{\Omega}^1_{V/Y})$ est localement libre de type fini, d'où 2.4.

\textbf{Corollaire 2.5.} \uline{Soit} $f:X\to Y$ \uline{un morphisme d'intersection complète lissifiable}, \uline{et} $f=g\circ i$ \uline{une factorisation de} $f$ \uline{par un} $Y$-\uline{schéma lisse} $V$. \uline{Soit} $\underline{J}$ \uline{l'idéal de} $\underline{O}_V$ \uline{définissant} $i$. \uline{Alors}, \uline{si} $K^{\bullet}(X)$ \uline{désigne le groupe de Grothendieck des} $\underline{O}_X$-\uline{Modules localement libres de type fini} (\uline{resp}. \uline{des complexes parfaits sur} $X$), \uline{l'élément} $T_f\in K^{\bullet}(X)$ \uline{défini par} :
\[
T_f=\operatorname{cl}(i^*(\underline{\Omega}^1_{V/Y}))-\operatorname{cl}(\underline{J}/\underline{J}^2)
\]
\[
(\text{resp. }T_f=\operatorname{cl}(\underline{L}^{X/Y}_\bullet))
\]
\uline{ne dépend pas de la factorisation de} $f$ \uline{choisie}. \uline{La valeur en un point} $x$ \uline{de} $X$ \uline{de l'augmentation de} $T_f$ \uline{n'est autre que la dimension relative virtuelle de} $X$ \uline{sur} $Y$ \uline{au point} $x$ (1.9).

L'élément $T_f$ est bien défini dans les deux cas envisagés d'après 2.4, et sa classe dans $K^{\bullet}(X)$ ne dépend pas de la factorisation, d'après 2.2, dans le cas respé, et d'après la démonstration de 2.2 dans le cas non respé. La dernière assertion de 2.5 est triviale.

\textbf{Proposition 2.6.} \uline{Soient} $f:X\to Y$, \uline{et} $g:Y\to Z$ \uline{deux morphismes d'intersection complète lissifiables}, \uline{et supposons de plus la condition suivante remplie}:

